\section{Related Work}
\label{sec:related-work}

\subsection{Tabular Foundation Models and In-Context Learning}

TabPFN~\cite{hollmann2023tabpfn} introduced a transformer pretrained once on
millions of synthetic datasets that performs supervised tabular classification
through in-context learning, the training set is presented as a prompt, and a
single forward pass yields predictions for the entire test set, without any
gradient-based fitting on the target data. The follow-up TabPFN~v2
release~\cite{hollmann2025tabpfnv2} scaled the pretraining prior and
architecture, reporting accuracy that matches or exceeds tuned gradient-boosted
ensembles on datasets of up to ten thousand rows while running orders of
magnitude faster. A concurrent empirical comparison across classical,
deep-learning, and tabular foundation models on imbalanced clinical
data~\cite{brima2026robustness} found that TabPFN~v2 and a related
in-context model, TabICL, rank competitively with tuned gradient-boosted
ensembles on one imbalanced clinical dataset but fall to intermediate,
less stable ranks on another, with XGBoost the only model family stable
across both. That comparison, like the TabPFN~v2 release itself, evaluates
imbalance robustness only under clean, uncorrupted labels, leaving open
whether either family's apparent robustness to imbalance survives once
label quality is also compromised, the question this paper asks directly.
Because a Prior-Data Fitted Network's predictions are shaped
entirely by the class distribution implicit in its pretraining prior together
with the context it is given, its behavior under a training-time class
distribution that differs from what the prior assumes is a distinct concern
from that of a conventionally trained classifier, and has only recently begun
to receive dedicated attention. Hu and Ghelichi~\cite{hu2026noiseimmunity}
take a complementary approach, analyzing TabPFN's internal attention
mechanisms directly under synthetic label-quality degradation and finding
that its attention stays structured and its ROC-AUC stays high without any
correction applied at all. Their analysis targets the model's internals
rather than the correction methods practitioners actually deploy on top of
it, the gap this paper's audit addresses.

DistPFN~\cite{lee2026distpfn} is the newest instance of a much older idea,
adjusting a classifier's output probabilities to match a known or assumed
target class prior, without retraining, first formalized for conventional
probabilistic classifiers by Saerens et al.~\cite{saerens2002adjusting}. To
our knowledge, DistPFN is the first method to bring this idea to an
in-context tabular foundation model, where the ``prior'' being adjusted for
is implicit in the pretrained transformer rather than an explicit model
parameter: it rescales TabPFN's predicted class probabilities to counteract
label shift between the training context and the deployment distribution,
with a temperature-scaled variant (DistPFN-T) that adapts the adjustment
strength to the discrepancy between the observed and reference priors. Both
variants require no retraining or architectural change, which places them in
the same class-imbalance setting that classical resampling and
cost-sensitive techniques were designed for on conventional classifiers,
despite operating through a completely different mechanism.

\subsection{Correcting Class Imbalance}

The dominant strategies for class imbalance fall into three families:
data-level resampling, algorithm-level cost adjustment, and, for
in-context tabular models, the test-time posterior corrections described
above. SMOTE~\cite{chawla2002smote} synthesizes new minority-class examples by
interpolating between real minority points and their nearest minority
neighbors in feature space, rather than duplicating existing rows, and remains
one of the most widely adopted correction techniques nearly a quarter-century
after its introduction. Random oversampling, by contrast, corrects the
training distribution by duplicating existing minority rows without
synthesizing new feature vectors. Both are implemented as standard,
interchangeable components in imbalanced-learn~\cite{lemaitre2017imbalancedlearn},
a widely used toolbox implementing both corrections. Cost-sensitive
learning~\cite{elkan2001costsensitive} instead leaves the training
distribution untouched and reweights the loss so that minority-class errors
are penalized more heavily than majority-class errors in proportion to the
observed class ratio. XGBoost's \texttt{scale\_pos\_weight}
parameter~\cite{chen2016xgboost} is a direct, widely used implementation of
this idea. Whether these
corrections reliably help at all, even when every label is trustworthy, is
itself contested: Elor and Averbuch-Elor~\cite{elor2022smote} find that with
a sufficiently strong classifier and well-chosen hyperparameters, balancing
the training set often provides no benefit over the imbalanced data as-is,
and can be actively counterproductive, a finding that already complicates
the assumption that these corrections are safe by default, before label
noise is introduced at all. All of these methods share a further, implicit
assumption that is central to this paper, that the
observed class labels used to compute the resampling ratio, the interpolation
neighborhood, or the cost ratio are themselves trustworthy. These classical corrections were not originally designed for robustness
to label noise. Subsequent empirical work, discussed later in this
section, has evaluated several sampling approaches under noisy-label
settings, but considerably less is known about the newer test-time
posterior corrections built for tabular foundation models.

\subsection{Learning Under Label Noise}

A separate literature studies classification when the observed labels
themselves are corrupted, independently of any class imbalance. Early
theoretical work formalized class-conditional label noise, in which the
probability that a true label is flipped depends only on its true class, and
established conditions under which consistent classifiers can still be
recovered~\cite{natarajan2013noisylabels}. Asymmetric label noise, where the
flip probabilities differ across classes or flip in only one direction, was
shown to require different identifiability and denoising strategies than the
symmetric case~\cite{scott2013asymmetricnoise}, indicating that the direction
and structure of label corruption, not only its rate, can materially change
how a classifier should be corrected. Frénay and
Verleysen~\cite{frenay2014labelnoise} survey sources and taxonomies of label
noise and its effects on classifier robustness prior to the deep learning era,
while Song et al.~\cite{song2022noisylabelssurvey} survey the substantially
larger body of work on training deep networks to be robust to noisy labels,
including loss correction, sample selection, and label correction approaches.
Both surveys treat label noise as a phenomenon to be corrected or tolerated
during training. Neither the surveys nor the methods they catalog examine how
label noise interacts with a class-imbalance correction applied
\emph{alongside} it, which is the specific interaction this paper studies,
not by proposing a new noise-robust training method, but by auditing whether
existing, widely deployed imbalance corrections remain net-beneficial once
that assumption is violated.

\subsection{The Understudied Intersection: Imbalance Correction Under Noisy Labels}

Empirically auditing whether imbalance corrections hold up under label
noise has important precedent. Van Hulse and
Khoshgoftaar~\cite{vanhulse2009knowledge} trained more than two million
classifiers across seven datasets, eleven learners, and seven sampling
techniques, including random oversampling and SMOTE, while varying the
location and level of class noise, showing that the class receiving the
noise materially affects how learning and sampling procedures behave and
establishing that imbalance mitigation can interact strongly with label
noise in classical learning settings. Our study addresses a
complementary setting, We control the total training-set size across
imbalance ratios so that changes in training-set size are not confounded
with changes in class imbalance, and we extend the analysis to tabular
foundation models and recently proposed test-time posterior corrections
such as DistPFN and DistPFN-T.

A smaller body of work has begun to treat class imbalance and label noise as
jointly occurring problems rather than independent ones. Koziarski et
al.~\cite{koziarski2020ccr} combine a resampling step with a label-cleaning step
specifically for multi-class imbalanced data known to contain label noise.
Chen et al.~\cite{chen2021rsmote} and Li et al.~\cite{li2024wrnd} take a
related but more direct approach, redesigning SMOTE itself so that its
neighborhood-based interpolation is weighted by local sample density and
downweights or filters out likely-noisy minority points, rather than pairing
an unmodified SMOTE with a separate cleaning stage. Khanal et
al.~\cite{khanal2024activelabel} combine active learning with
noisy-label training to handle both problems together in imbalanced medical
image classification. Liu et al.~\cite{liu2024imbalancednoisy} observe that
standard loss-based sample selection for noisy-label training is itself biased
against minority classes, since under-learned tail classes and noise-corrupted
classes are both associated with high training loss and are consequently
difficult to tell apart. They propose a class-balanced sample selection
procedure with confidence-based augmentation to correct this bias during
training. A recent line of work inverts the relationship between the two
problems entirely: Hu et al.~\cite{hu2025beneficialnoise} show that
\emph{deliberately} injecting an asymmetric label noise model, engineered
specifically to counteract a classifier's biased decision boundary, can
rebalance an imbalanced classifier without resampling at all. Their result
and ours are not in tension. They show that noise designed in advance to
help, with its direction and magnitude chosen to correct a known bias, can
substitute for a resampling or reweighting step; we show that the standard
resampling and reweighting steps already in deployment cannot tell that
engineered noise apart from the generic, uncontrolled label errors a real
dataset actually contains, and so extend the same corruption-tolerance
assumption to the corrections themselves rather than to the raw
classifier. Each of these lines of work proposes a new algorithm for the
training-time mechanism of handling noisy labels under imbalance. None of them
evaluates whether the standard \emph{imbalance corrections already in widespread use}
(SMOTE, random oversampling, cost-sensitive reweighting, or, for the newer
class of tabular foundation models, test-time posterior adjustment and
class-balanced downsampling) continue to help, are neutral, or actively harm
performance once a controlled fraction of the labels they operate on are
themselves wrong. What remains open, to our knowledge, is whether these
corrections, applied to tabular foundation models under a directed,
asymmetric label-error rate, still improve on doing nothing at all rather
than merely being neutral or harmful. Closing that gap is what this
paper's controlled, cross-method audit is designed to do.
